\documentclass{mnnit}

\usepackage{graphicx}
\usepackage{booktabs}
\usepackage{amsmath}
\usepackage{array}
\usepackage{url}
\usepackage{etoolbox}
% Single spacing inside figures and tables (the class used to do this; see mnnit.cls).
\AtBeginEnvironment{figure}{\def\baselinestretch{1}\normalsize}
\AtBeginEnvironment{table}{\def\baselinestretch{1}}

\graphicspath{{figures/}}
\setlength{\emergencystretch}{2.5em}

% ---------------------------------------------------------------------------
% Group-project front matter. mnnit.cls writes the undertaking in the first
% person singular and puts \@author inside a sentence, so both pages are
% redefined here for four authors.
% ---------------------------------------------------------------------------
\newcommand{\authorsinline}{Shaurya Pratap Singh (20233262), Surya Pratap Singh (20233283), Vishesh Vigh (20233583), and Sehar Majid Abdi (20233561)}

\makeatletter
\def\undertakingpage{\newpage%
    \renewcommand{\baselinestretch}{1.3}
    \vspace*{\fill}
    \centerline{\LARGE\bf UNDERTAKING}
    \vspace{2cm}
    \begin{center}
    \begin{minipage}{5.2in}
    \parindent=0pt
    We declare that the work presented in this report titled
    ``{\sl\@title}'', submitted to the \mbox{\@dept} Department, \@college, for
    the award of the {\bf\emph{\@degree}} degree in \mbox{\bf\emph{\@specialization}},
    is our original work. We have not plagiarized or submitted the same work for
    the award of any other degree. In case this undertaking is found incorrect, we accept
    that our degree may be unconditionally withdrawn.\par
    \vspace*{1.5cm}
    \@submitdate\\Prayagraj\par
    \vspace*{1.8cm}
    \begin{tabular}{@{}p{2.5in}p{2.5in}@{}}
    \rule{2.1in}{0.4pt} & \rule{2.1in}{0.4pt}\\
    Shaurya Pratap Singh (20233262) & Surya Pratap Singh (20233283)\\[1.6cm]
    \rule{2.1in}{0.4pt} & \rule{2.1in}{0.4pt}\\
    Vishesh Vigh (20233583) & Sehar Majid Abdi (20233561)
    \end{tabular}
    \end{minipage}
    \end{center}
    \vspace*{\fill}
}

\long\def\signature#1{%
\begin{center}
\begin{minipage}{5.2in}
\parindent=0pt
Certified that the work contained in the report titled
``{\sl\@title}'', by \authorsinline, has been carried out under my
supervision and that this work has not been submitted elsewhere for a
degree.
\par
\vspace{3cm}
\hbox to 5.2in{\@submitdate\hfill
    \vbox{\hsize 3in\vrule width 2in height0.4pt\par#1}\hss}
\end{minipage}
\end{center}
}
\makeatother

% Citations. The class redefines thebibliography after loading natbib, which
% leaves \cite unresolved (the template's own PDF prints "?"). References are
% therefore a numbered list cited with \rc{key}, which needs no BibTeX run.
\makeatletter
\newcommand{\rc}[1]{[\def\rc@sep{}\@for\rc@key:=#1\do{\rc@sep\ref{ref:\rc@key}\def\rc@sep{, }}]}
\makeatother

\begin{document}
\title{Neuro-Symbolic Reasoning over Patient-Specific Knowledge Graphs for Breast Cancer Progression Prediction}
\author{\begin{tabular}{ll}
Shaurya Pratap Singh & (20233262)\\
Surya Pratap Singh & (20233283)\\
Vishesh Vigh & (20233583)\\
Sehar Majid Abdi & (20233561)
\end{tabular}}
\supervisor{Dr. Shashwati Banerjea}
\specialization{Computer Science \& Engineering}

\figurespagetrue
\tablespagetrue
\beforepreface

\prefacesection{Abstract}
Predicting whether and when a breast cancer will progress after treatment is a central question in oncology,
and a prediction is only useful to a clinician if it can be explained. This project builds patient-specific
knowledge graphs for the 1,083 breast invasive carcinoma patients of the TCGA PanCancer Atlas 2018 cohort and
fuses each one with PrimeKG, a biomedical knowledge graph of 8.1 million relations. Every patient graph links
the patient to their tumour sample and to their mutated and copy-number-altered genes from a 500-gene panel.
These genes are joined to PrimeKG pathways, drugs, and breast-disease entities through their NCBI Entrez
identifiers.

On these fused graphs we benchmark seven graph neural network architectures (GCN, GraphSAGE, GATv2, R-GCN, RGAT,
GIN, and a graph transformer) under three objectives (binary cross-entropy, a differentiable concordance index,
and the Cox partial likelihood). We also add patient-conditioned graph structure learning. Discrimination
saturates near a concordance index of 0.67 across every model family, and the clinical variables carry most of
the signal.

We therefore turn the system into a neuro-symbolic one. A survival GNN supplies neural edge importances. For the
symbolic part we go beyond IF--THEN rule induction and compare four survival-aware reasoners over 422 first-order
ground atoms, including relational PrimeKG atoms such as $\mathit{Altered}(p,g) \wedge \mathit{InPathway}(g,w)$:
a Markov logic network whose formula weights are learned with the Cox partial likelihood, probabilistic soft logic
with a collective patient-similarity rule, a policy-gradient rule generator with a replication reward, and a
hybrid in which reinforcement learning proposes the Markov logic formulas. The hybrid is the rule engine. A fused
Cox model, chosen by the one-standard-error rule, produces each patient's risk and survival curve. Accelerated
failure time models fitted on progressors estimate the month of progression. Each explanation path is scored as
0.6 times its neural importance plus 0.4 times the weights of the learned formulas it grounds, and a prominence
cutoff selects the top-K paths.

The system writes a hospital-style PDF report for every patient, grouped by recorded progression timeline. All
predictions are out-of-fold. The Markov logic engines reach a concordance of 0.63 against 0.62 for RuleKit, with
69\% of their formulas reasoning over PrimeKG relations against 10\% of RuleKit's rules. The fused model matches the
best clinical model (concordance 0.664 against 0.675, within one standard error). The progression-month estimate
is within 12 months of the recorded month for 48\% of progressors, against 37\% for a constant baseline. Every
report grounds its risk estimate in explicit knowledge-graph paths and weighted logic formulas.

\prefacesection{Acknowledgements}
We express our sincere gratitude to our supervisor, Dr.\ Shashwati Banerjea, for guidance, patience with
our many iterations, and above all for steering the project from raw model accuracy towards neuro-symbolic
reasoning and interpretable patient reports. This feedback shaped every milestone of this work.

We thank the Department of Computer Science and Engineering, Motilal Nehru National Institute of Technology
Allahabad, for the academic environment and resources that made this project possible. We are grateful to the
TCGA Research Network and the cBioPortal team for making the BRCA PanCancer Atlas data public, to the Zitnik Lab
at Harvard for PrimeKG, and to the developers of PyTorch Geometric, lifelines, scikit-survival, and RuleKit, on
whose open-source work this project is built.

\afterpreface

% ===========================================================================
\chapter{Introduction}
% ===========================================================================

Breast cancer is the most commonly diagnosed cancer worldwide, with an estimated 2.3 million new cases in 2020
\rc{globocan}. Its course after primary treatment varies widely. Some patients never relapse, some progress within
months, and some progress a decade later. Clinicians estimate this risk from pathological stage, lymph-node
status, and molecular subtype \rc{pam50}. Genome-scale profiling, however, records hundreds of somatic alterations
per tumour, and most of them have no established prognostic meaning on their own.

Biomedical knowledge graphs offer a way to read those alterations in context. PrimeKG \rc{primekg} integrates
twenty curated resources into a single graph of 129,375 nodes and over 8 million relations. It connects genes and
proteins to pathways, drugs, diseases, and phenotypes. An altered gene in a tumour can therefore be connected to the
pathways it participates in, to the drugs that target it, and to the breast diseases it has been linked to. Graph
neural networks (GNNs) \rc{gcn,graphsage} can learn from such structure, and explainers such as GNNExplainer
\rc{gnnexplainer} can say which parts of a graph drove a prediction.

\section{Motivation}
A prognostic model is only useful at the bedside if a clinician can see why it reached its conclusion.
Purely neural models over large graphs give a score without reasons. Purely symbolic models, such as survival
rules, are readable but cannot weigh thousands of interacting facts. Neuro-symbolic AI \rc{neurosymbolic} combines
the two. Neural components learn from data, and symbolic components state knowledge explicitly and can be checked.

This project started as a pure GNN benchmark. After several iterations, including relation-aware architectures,
concordance training, and graph structure learning, its predictive performance saturated. Following our
supervisor's direction, the final phase builds a neuro-symbolic layer and turns its output into per-patient
reports. The reports explain every prediction with knowledge-graph paths and weighted logic formulas, and they group
patients by when they progressed.

\section{Problem Statement}
Given the clinical record and the somatic mutation and copy-number profile of a breast cancer patient, together
with the PrimeKG biomedical knowledge graph, the system must:
\begin{enumerate}
  \item estimate the patient's risk of disease progression and their progression-free survival curve;
  \item estimate the month in which progression occurs, given that it occurs;
  \item explain the estimate through explicit, verifiable knowledge-graph paths and weighted logic formulas; and
  \item present the result as a readable clinical report, grouped by progression timeline.
\end{enumerate}
Outcome variables must never leak into the model inputs. Every reported prediction must come from a model that
did not see that patient during training.

\section{Objectives}
\begin{itemize}
  \item Build validated patient-specific graphs from the cBioPortal TCGA BRCA PanCancer Atlas 2018 study \rc{cbioportal,cbioportal2}.
  \item Fuse each patient graph with a filtered PrimeKG neighbourhood through NCBI Entrez identifiers.
  \item Benchmark GNN architectures and training objectives on leakage-safe progression labels.
  \item Build a neuro-symbolic model that combines a survival GNN with survival-aware logical reasoning, comparing Markov logic \rc{mln}, probabilistic soft logic \rc{psl}, and reinforcement-learned rules against a RuleKit baseline \rc{rulekit}.
  \item Estimate the month of progression and evaluate it against the recorded month.
  \item Generate a hospital-style PDF report for every patient, binned by progression timeline.
\end{itemize}

\section{Contributions}
\begin{enumerate}
  \item A reproducible pipeline from raw cBioPortal tables to validated, leakage-checked patient graphs and PrimeKG-fused PyTorch Geometric graphs, with 58 unit tests.
  \item A benchmark of seven GNN architectures under three objectives, a horizon sweep, graph structure learning, and an external microarray benchmark, all with out-of-fold evaluation.
  \item A neuro-symbolic survival model. Its symbolic engine is a \emph{survival Markov logic network} whose
        first-order formulas, including multi-hop PrimeKG conjunctions, are proposed by a reinforcement-learning
        rule generator rewarded for replicated log-rank separation, and whose weights are learned with the Cox
        partial likelihood. It is compared with survival probabilistic soft logic with collective inference, RL
        rules alone, and RuleKit. Explanation paths are scored as 0.6 times their GNNExplainer importance plus 0.4
        times the weights of the formulas they ground.
  \item An honest model-selection procedure (the one-standard-error rule) and a progression-month estimator evaluated against recorded outcomes.
  \item 1,083 patient reports, group summaries, a cohort overview, and an executed notebook that documents the work.
\end{enumerate}

\section{Organisation of the Report}
Chapter~\ref{ch:related} reviews related work on cancer genomics data, knowledge graphs, graph neural networks,
survival analysis, and rule learning. Chapter~\ref{ch:proposed} describes the proposed system, from data to
reports. Chapter~\ref{ch:results} gives the experimental setup and analyses the results of every milestone.
Chapter~\ref{ch:conclusion} concludes and outlines future work. The appendix lists the commands that reproduce
every result.

% ===========================================================================
\chapter{Related Work}\label{ch:related}
% ===========================================================================

\section{Cancer Genomics Data}
The Cancer Genome Atlas (TCGA) PanCancer Atlas re-analysed more than 10,000 tumours from 33 cancer types with
uniform pipelines \rc{hoadley}. Liu et al.\ \rc{tcgacdr} curated its clinical data into a resource with four
standardised endpoints: overall survival, disease-specific survival, disease-free interval, and
progression-free interval. They recommended progression-free interval for breast cancer because of its event
count and follow-up. cBioPortal \rc{cbioportal,cbioportal2} distributes these data as tab-delimited study
bundles. Copy-number calls in these bundles are discretised by GISTIC2.0 \rc{gistic} into homozygous deletion,
hemizygous deletion, neutral, gain, and high-level amplification. The PAM50 intrinsic subtypes \rc{pam50}
(Luminal A, Luminal B, HER2-enriched, Basal-like, Normal-like) are the standard molecular classification of breast
cancer.

\section{Biomedical Knowledge Graphs}
PrimeKG \rc{primekg} is a precision-medicine knowledge graph built from 20 resources, including Reactome pathways
\rc{reactome}, DrugBank drug targets \rc{drugbank}, and curated disease-gene associations. It describes 17,080
diseases with 4,050,249 relations. Its nodes keep source identifiers, such as NCBI Entrez for genes, so patient data
can be aligned to it without fuzzy name matching. Knowledge graphs supply the symbolic, human-readable layer that
neuro-symbolic systems reason over.

\section{Graph Neural Networks}
Graph convolutional networks (GCN) \rc{gcn} propagate degree-normalised neighbour features. GraphSAGE \rc{graphsage}
learns separate self and neighbour transforms and generalises inductively to unseen graphs, which suits a setting
with one graph per patient. Graph attention networks learn a weight per edge, and GATv2 \rc{gatv2} fixes the
static-attention limitation of the original formulation. Relational GCN (R-GCN) \rc{rgcn} and relational graph
attention (RGAT) \rc{rgat} give each relation type its own weights, which is the classical treatment of
heterogeneous knowledge graphs. The graph isomorphism network (GIN) \rc{gin} uses sum aggregation and is as
expressive as the Weisfeiler--Lehman test. TransformerConv \rc{transformerconv} brings multi-head dot-product
attention with edge features to message passing. All of these are implemented in PyTorch Geometric \rc{pyg}.
GNNExplainer \rc{gnnexplainer} learns a soft mask over edges that preserves a model's prediction, which yields an
importance score per edge.

Graph structure learning \rc{gslsurvey,lds} learns which edges a GNN should use, instead of taking the input
graph as given. This is relevant when a knowledge-graph neighbourhood may be too large or too noisy.

\section{Survival Analysis}
Progression is a time-to-event outcome with right censoring: a patient without a recorded progression may still
progress after their last follow-up. The Kaplan--Meier estimator \rc{km} estimates a survival curve non-parametrically.
The log-rank test \rc{logrank} compares the curves of two groups. The Cox proportional hazards model \rc{cox}
relates covariates to a hazard ratio and is fitted by maximising the partial likelihood. Accelerated failure time
(AFT) models \rc{aft} instead model the logarithm of the event time directly. Random survival forests \rc{rsf}
extend random forests to censored data. DeepSurv \rc{deepsurv} trains a neural network on the Cox partial
likelihood, the approach our survival GNN follows.

Discrimination is measured by Harrell's concordance index \rc{harrell}: among comparable pairs of patients, the
fraction in which the one who progressed earlier received the higher risk. Uno's C \rc{uno} reweights pairs by the
inverse probability of censoring, which removes the dependence on the censoring distribution. Python
implementations are provided by lifelines \rc{lifelines} and scikit-survival \rc{sksurv}.

\section{Rule Learning for Survival Data}
LR-Rules \rc{lrrules} is a sequential-covering algorithm for survival data. It grows each rule by adding the
condition that maximises the log-rank separation between covered and uncovered patients. Each rule's conclusion is
the Kaplan--Meier curve of the patients it covers. RuleKit \rc{rulekit} packages LR-Rules, together with
classification and regression rule learners, as a Java library with a Python interface. Survival rules are
readable statements such as ``IF nodal stage is N3 AND aneuploidy is below 26.5 THEN the progression-free survival
curve is \ldots'', which makes them a natural symbolic component.

\section{Statistical Relational Learning: Markov Logic and Soft Logic}
Markov logic networks (MLNs) \rc{mln} attach a real weight $w_k$ to each first-order formula $F_k$ and define a
log-linear distribution over possible worlds, $P(x) \propto \exp\big(\sum_k w_k\, n_k(x)\big)$, where $n_k(x)$ is
the number of true groundings of $F_k$. A formula with a large weight is a soft constraint: worlds that violate it
are less probable but not impossible. MLN learning has two parts: structure learning finds the formulas
\rc{mlnstructure}, and weight learning fits $w$, usually by maximising a pseudo-likelihood.

Probabilistic soft logic (PSL) \rc{psl} relaxes truth values to $[0,1]$ with the {\L}ukasiewicz t-norm. A weighted
rule $w: B \Rightarrow H$ becomes the hinge potential $w \max(0, B - H)^q$, its distance to satisfaction. The resulting
hinge-loss Markov random field has convex MAP inference, so PSL scales to collective reasoning over many
interrelated entities, for example propagating a label over a similarity graph.

\section{Reinforcement Learning for Rules and Paths}
Policy-gradient reinforcement learning \rc{reinforce} optimises a stochastic policy from scalar rewards without a
differentiable objective. In knowledge-graph reasoning, DeepPath \rc{deeppath} and MINERVA \rc{minerva} train
agents to walk relation paths, and RNNLogic \rc{rnnlogic} trains a recurrent rule generator jointly with a
reasoning predictor that scores its rules. We follow the generator idea with a gated recurrent unit \rc{gru} that
writes conjunctive survival rules, rewarded by log-rank separation, and we use its output as the structure learner
of a survival MLN.

\section{Neuro-Symbolic AI}
Garcez and Lamb \rc{neurosymbolic} describe neuro-symbolic AI as the integration of neural learning with symbolic
knowledge representation and reasoning, so that systems can both learn from data and explain themselves with
explicit knowledge. In biomedicine this typically means grounding neural predictions in curated knowledge graphs
and rules. Our design follows that pattern. The neural component scores knowledge-graph paths, the symbolic
component supplies population evidence for the facts on those paths, and a survival rule learner states the
cohort-level regularities explicitly.

\section{Gene-Expression Prognosis Cohorts}
GSE2034 \rc{gse2034} (286 lymph-node-negative patients, Affymetrix arrays) and GSE2990 \rc{gse2990} (189 patients
with relapse-free survival) are classic external breast cancer prognosis cohorts. We use them as an external
benchmark for the graph models.

% ===========================================================================
\chapter{Proposed Work}\label{ch:proposed}
% ===========================================================================

Figure~\ref{fig:pipeline} gives an overview of the system. Patient graphs are built from TCGA and fused with a
filtered PrimeKG neighbourhood. A survival GNN and a survival rule learner run on the same cohort and the same
folds. Their risks are fused into one Cox model, and every patient receives a report that combines the
prediction with GNN-scored knowledge-graph paths and the rules that fire for them.

\begin{figure}[htbp]
  \centering
  \includegraphics[width=\textwidth]{fig_pipeline.png}
  \caption{Overview of the neuro-symbolic progression system.}
  \label{fig:pipeline}
\end{figure}

\section{Data}
The study is the cBioPortal \texttt{brca\_tcga\_pan\_can\_atlas\_2018} bundle. We verified every data file
against the official SHA-256 checksums. Four tables are used:
\begin{itemize}
  \item \texttt{data\_clinical\_patient.txt}: 1,084 patients, 38 columns, including the PFS status and months;
  \item \texttt{data\_clinical\_sample.txt}: 1,084 primary tumour samples, 19 columns;
  \item \texttt{data\_mutations.txt}: 126,252 MAF mutation records with 111 annotation columns;
  \item \texttt{data\_cna.txt}: discrete GISTIC calls for 25,128 genes across 1,070 samples.
\end{itemize}
A streamed audit of all 23 tables in the bundle records their schemas and missingness. The expression,
methylation, protein, and structural-variant modalities were inspected and deliberately excluded to keep the
graph interpretable.

\paragraph{Leakage policy.} Outcome columns (OS, DSS, DFS, and PFS status and months, the new-tumour-event flag,
and the cancer status at last contact) are never features. A guard function rejects any feature whose name
matches an outcome column, and tests assert it.

\section{Patient Graph Construction}
Each patient is a \texttt{networkx} \texttt{MultiDiGraph} with the schema
Patient $\rightarrow$ Sample $\rightarrow$ Gene (Figure~\ref{fig:patientgraph}):
\begin{itemize}
  \item \textbf{Patient} nodes carry age, sex, PAM50 subtype, AJCC stage, and pathological M stage.
  \item \textbf{Sample} nodes carry the tumour type, OncoTree code, aneuploidy score, and a raw MAF record count.
  \item \textbf{Gene} nodes are keyed by HGNC symbol and Entrez identifier. The identifier pair avoids merging the
        duplicated symbol rows of the CNA matrix.
  \item \texttt{has\_mutation} edges preserve each MAF record (variant class, protein change, coordinates,
        alleles, caller filters). \texttt{has\_CNA} edges carry the GISTIC value.
\end{itemize}
Exact duplicate CNA calls are deduplicated. Contradictory calls for the same gene and sample are withheld and
logged rather than resolved arbitrarily. Every graph is validated before export.

\begin{figure}[htbp]
  \centering
  \includegraphics[width=0.85\textwidth]{fig_patient_graph.png}
  \caption{A Milestone~1 patient graph (TCGA-A2-A0SW): one patient, one tumour sample, and its mutated genes.}
  \label{fig:patientgraph}
\end{figure}

\section{Gene Panel}
To bound graph size, genes are restricted to a 500-gene panel ranked by an integrated alteration score: the number
of samples with a mutation plus the number with an extreme CNA ($-2$ or $+2$). Twenty-seven literature breast
cancer drivers are always included: TP53, PIK3CA, GATA3, CDH1, ERBB2, PTEN, MAP3K1, BRCA1, BRCA2, RB1, CCND1, MYC,
ATM, ESR1, AKT1, MAP2K4, KMT2C, ARID1A, NF1, EGFR, FGFR1, FOXA1, CBFB, RUNX1, PIK3R1, SF3B1, and BAP1. For
modelling, CNA evidence is limited to homozygous deletion and high-level amplification.

\section{Fusion with PrimeKG}
PrimeKG's 8,100,498 edges are streamed once and filtered to a 4,570-edge panel slice that keeps four relations:
\begin{itemize}
  \item \texttt{protein\_protein} edges between two panel genes (772);
  \item \texttt{pathway\_protein} edges from panel genes to Reactome pathways (2,410);
  \item \texttt{drug\_protein} edges from panel genes to DrugBank drugs (858);
  \item \texttt{disease\_protein} edges to diseases whose name contains ``breast'' or ``mammary'' (530).
\end{itemize}
Patient genes join PrimeKG proteins on Entrez identifier. The fused graph (Figure~\ref{fig:fused}) adds, for each
altered gene, its panel protein partners, its pathways, drugs, and breast-disease links, and the unaltered panel
genes that share a pathway with it (two hops). Every relation is also stored in reverse so that three
message-passing layers can carry pathway context back to the patient node. Over the cohort the median fused graph
has 632 nodes and 2,526 directed edges.

Node features have one width: a node-type one-hot, a 19-field clinical block on the patient node, a five-field
alteration block on gene nodes (mutated, amplified, deleted, $\log(1+\text{mutation count})$, known driver), and a
trainable embedding for every gene, pathway, drug, and disease.

\begin{figure}[htbp]
  \centering
  \includegraphics[width=0.9\textwidth]{fig_fused_graph.png}
  \caption{Part of a fused graph: altered genes of one patient connected to PrimeKG partner proteins, pathways, and a breast-disease node.}
  \label{fig:fused}
\end{figure}

\section{Progression Labels}
The target is the progression-free interval (\texttt{PFS\_STATUS}, \texttt{PFS\_MONTHS}). 1,083 patients have a
usable status, and 145 of them have a recorded progression. A censored patient was not seen to progress during
follow-up. Treating censoring as ``never progresses'' biases a binary model, so two endpoints are used:
\begin{itemize}
  \item \textbf{36-month landmark} (\texttt{pfs\_36m}): label 1 for progression at or before 36 months and 0 for
        follow-up beyond 36 months. Patients censored before 36 months are dropped, which leaves 498 patients and
        89 events.
  \item \textbf{Full follow-up} (\texttt{pfs}): the time and event pair for all 1,083 patients, used by the survival
        objectives and by the neuro-symbolic model.
\end{itemize}
The PFS months define the target and the comparable pairs. They are never model inputs.

\section{Graph Neural Network Models}
Seven architectures share one design: three message-passing layers of width 64, dropout 0.3, and a readout that
concatenates the patient node with mean- and max-pooling over the patient and altered genes. A linear clinical head
is added to the graph score through a gate that starts near zero, so training begins from the clinical score. The
architectures are GCN, GraphSAGE, GATv2 (four heads), R-GCN and RGAT over the eight relation types, GIN (a control
that ignores relations), and a graph transformer that reads a learned relation embedding as an edge feature. A
multilayer perceptron over the clinical block and a 500-gene alteration multi-hot, and a clinical logistic
regression, are the non-graph baselines.

\subsection{Training Objectives}
\begin{itemize}
  \item \textbf{Weighted binary cross-entropy} on the landmark label.
  \item \textbf{Differentiable C-index}: for comparable pairs $(i,j)$ with $t_i < t_j$ and $\delta_i = 1$, minimise
        $1 - \frac{1}{|P|}\sum_{(i,j)\in P}\sigma\!\left(\frac{r_i - r_j}{\tau}\right)$ over each mini-batch.
  \item \textbf{Cox partial likelihood} (Breslow ties):
  \begin{equation}
    \mathcal{L} = -\frac{1}{N_e}\sum_{i:\,\delta_i=1}\Big(r_i - \log\!\!\sum_{j:\,t_j \ge t_i}\! e^{r_j}\Big),
  \end{equation}
  where $r_i$ is the model's log-risk for patient $i$.
\end{itemize}

\section{Patient-Conditioned Graph Structure Learning}
A fixed two-hop neighbourhood may include PrimeKG edges that are irrelevant to a given patient. The
graph-structure-learning (GSL) variant retrieves up to three hops and learns a weight for every PrimeKG edge from
the edge's endpoints, its relation, and the patient's clinical context. In the hop-prior version, edges within the
fixed two hops start at weight 1, so training starts from the fixed graph. The third hop sits behind a
patient-specific gate that starts closed (about 0.12) and must cross 0.5 to contribute. The learned weight
multiplies messages in GCN, GraphSAGE, and GATv2. The relational convolutions in PyTorch Geometric accept no edge
weight, so GSL with them is refused rather than silently ignored.

\section{Neuro-Symbolic Progression Model}
The final system has four stages, all run over one cohort (1,083 patients) and one set of stratified five-fold
splits. Every prediction, rule, and path score for a patient therefore comes from models trained on the four
folds that exclude that patient.

\subsection{Neural Component: Survival GNN}
GCN, GraphSAGE, and R-GCN are trained on the fused graphs with the Cox objective, and GraphSAGE also with the
C-index objective, on the full follow-up. Mini-batches of 128 graphs are used for the Cox loss so that each batch
holds enough events for its risk sets. Early stopping uses the validation C-index. The best model by fold-mean
C-index supplies the neural risk. For every patient, GNNExplainer \rc{gnnexplainer}, run on the fold model that
held that patient out, supplies a soft importance for every edge of the patient's graph.

\subsection{Symbolic Baseline: Predicates and RuleKit Rules}
Each patient is described by 188 readable predicates:
\begin{itemize}
  \item \textbf{Clinical (10)}: age, T stage, N stage, M stage, AJCC stage, PAM50 subtype, histology (ductal,
        lobular, other), radiation therapy, aneuploidy score, and tumour mutational burden.
  \item \textbf{Genomic (58)}: mutations of the 27 drivers, driver copy-number changes, amplicons, and the number of
        altered panel genes. Co-amplified genes whose patient sets overlap with Jaccard index $\ge 0.8$ are merged
        into one predicate named after a driver, for example ``MYC amplified (+27 co-altered genes)''. Mutations
        of non-driver genes, mostly long passenger genes such as TTN or RYR2, are excluded.
  \item \textbf{Knowledge-graph (120)}: ``an altered gene of this patient is annotated to Reactome pathway $P$''
        (116) or ``\ldots is linked to breast disease $D$'' (4). Predicates are kept when 5--60\% of patients carry
        them, and near-duplicates are merged.
\end{itemize}
Predicate selection uses prevalence only, never the outcome. Binary predicates enter RuleKit as ``yes'' or
missing, so a rule can say ``TP53 mutated'' but cannot be built from the absence of a rare mutation.

As a baseline, RuleKit's LR-Rules \rc{lrrules,rulekit} is fitted on each training fold (minimum coverage 5\% of new patients, at
most three conditions per rule). A new patient's prediction is the average Kaplan--Meier curve of the rules that
cover them. The rule risk is the negative restricted mean survival time up to 120 months:
\begin{equation}
  r^{\text{rule}}_i = -\int_0^{120} \hat{S}^{\text{rule}}_i(t)\,dt .
\end{equation}

\subsection{Knowledge-Graph Fact Evidence}\label{sec:kgevidence}
For every PrimeKG pathway, drug, and disease node $v$, the fact $f_v$ is ``the patient has an altered gene directly
annotated to $v$''. On each training fold a vectorised two-group log-rank test compares carriers with
non-carriers, giving an observed/expected hazard-ratio estimate and a z statistic $z_v$. Facts with fewer than 15
carriers are skipped. In total 1,053 facts receive evidence.

\subsection{Logic Reasoners: Beyond IF--THEN Rules}\label{sec:logic}
RuleKit's crisp rules cover few patients, carry no weights, and rarely use the knowledge graph. We therefore build
four survival-aware reasoners over a shared set of first-order \emph{ground atoms} and compare them on the same
folds.

\paragraph{Ground atoms.} Every atom is a closed formula about a patient $p$, and its value $n(p)$ is its number of
true groundings. There are 26 clinical atoms such as $\mathit{NStage}(p, \mathrm{N3})$ or
$\mathit{Age}(p) \ge 65$, 57 genomic atoms such as $\mathit{Mutated}(p, \mathrm{TP53})$, and 339 relational PrimeKG
atoms
\begin{equation}
  n_w(p) = \big|\{\, g : \mathit{Altered}(p,g) \wedge \mathit{InPathway}(g,w) \,\}\big|,
\end{equation}
and likewise with $\mathit{TargetedBy}(g,d)$ for drugs and $\mathit{LinkedTo}(g,d)$ for breast diseases. Atoms
are kept at 3--90\% prevalence, which never uses the outcome. Atoms with identical carrier sets are merged, and
the merged atom keeps every PrimeKG node it stands for.

\paragraph{Survival Markov logic network (Cox-MLN).} Every formula concludes progression,
$w_k : \mathit{Body}_k(p) \Rightarrow \mathit{Progress}(p)$. The bodies are every atom plus up to 150
clinical/genomic $\times$ PrimeKG conjunctions mined on the training fold. With fully observed evidence and one
query atom per patient, MLN inference factorises patient by patient into the score
$\eta(p) = \sum_k w_k \log\!\big(1 + n_k(p)\big)$, where $\log(1+n)$ dampens patients with many co-amplified genes in
one pathway. The usual conditional likelihood would treat censored patients as non-progressors. We instead learn the
weights by maximising the elastic-net-penalised Cox partial likelihood \rc{elasticnet,coxnet}
\begin{equation}
  \max_{w}\; \sum_{i:\,\delta_i = 1}\Big(\eta_i - \log\!\!\sum_{j:\,t_j \ge t_i}\! e^{\eta_j}\Big)
  \;-\; \lambda\Big(\alpha\|w\|_1 + \tfrac{1-\alpha}{2}\|w\|_2^2\Big),
\end{equation}
with $\alpha = 0.5$ and $\lambda$ chosen by inner three-fold C-index. The $\ell_1$ term performs structure
learning: formulas whose weight reaches zero are pruned.

\paragraph{Survival probabilistic soft logic (S-PSL).} The query atom $\mathit{HighRisk}(p) = h_p \in [0,1]$ is
inferred from the program
\begin{align}
  w^+_k &: \mathit{Fact}_k(p) \Rightarrow \mathit{HighRisk}(p), \nonumber\\
  w^-_k &: \mathit{Fact}_k(p) \Rightarrow \neg\mathit{HighRisk}(p), \nonumber\\
  s^+ &: \mathit{Similar}(p,q) \wedge \mathit{HighRisk}(q) \Rightarrow \mathit{HighRisk}(p), \nonumber\\
  s^- &: \mathit{Similar}(p,q) \wedge \neg\mathit{HighRisk}(q) \Rightarrow \neg\mathit{HighRisk}(p),
\end{align}
plus a prior $\neg\mathit{HighRisk}(p)$. Facts are the 40 strongest atoms on the training fold, with PrimeKG
counts softened to $\min(1, n/2)$. $\mathit{Similar}(p,q)$ links each patient to their ten most similar training
patients by cosine similarity of atoms. For those neighbours $h_q$ is observed: it is a censoring-aware soft target,
the Kaplan--Meier probability of progression by 36 months given the patient's follow-up. MAP inference minimises
the squared-hinge energy by 40 projected-gradient steps. Inference is unrolled, so the rule weights are learned end
to end by soft-label cross-entropy plus a Cox term.

\paragraph{Reinforcement-learned rules (RL-Rules).} A GRU policy \rc{gru} writes a rule one atom at a time (up to
three atoms, or STOP), masking atoms of an already chosen clinical variable. It is trained with REINFORCE
\rc{reinforce}, a moving-average baseline, and an entropy bonus over 160 batches of 256 episodes per fold. The
training fold is split into a 70\% search part, halved into $A$ and $B$, and a 30\% validation part. The reward is
a \emph{stability reward}:
\begin{equation}
  R(\text{rule}) = \mathbb{1}\big[\operatorname{sign} z_A = \operatorname{sign} z_B\big]\,
  \frac{\min\!\big(|z_A|, |z_B|, 5\big)}{5} \;-\; 0.03\,(\text{length} - 1),
\end{equation}
where $z_A$ and $z_B$ are the log-rank statistics of the covered patients against the rest on the two halves.
Rules covering fewer than 4\% or more than 90\% of the fold earn nothing. Because tens of thousands of distinct rules
are sampled, the best training rewards are a multiple-testing artefact. Only rules whose separation replicates on the
untouched 30\% (same sign, $|z| \ge 1$) are kept, and a ridge Cox model on their activations gives the risk.

\paragraph{RL-structured survival MLN (engine).} The RL search proposes up to 150 replicated multi-atom formulas.
They are added to the Cox-MLN's candidates, and the Cox-MLN learns every weight. Reinforcement learning does the MLN's
structure learning and the survival likelihood does its weight learning, combining two of the three families.

\paragraph{Choosing the engine.} The reasoners are compared by fold-mean C-index and the one-standard-error rule
\rc{cart}. Among reasoners within one standard error of the best, the earliest in the preference order RL-structured
MLN $>$ MLN $>$ PSL $>$ RL rules $>$ RuleKit is the engine, which puts the richest reasoning first.

\subsection{Fusion and Model Selection}
The fused model is a ridge-penalised Cox model on the 19-field clinical block, the out-of-fold GNN risk, and the
out-of-fold risk of the logic engine. This is cross-validated stacking \rc{stacking}. It is compared with clinical Cox models (19
fields and an extended 27-column design), the neural and symbolic risks alone, and a random survival forest
\rc{rsf}. The report model is chosen by the one-standard-error rule \rc{cart}. Let $\bar{C}^*$ be the best
fold-mean C-index and $s^*$ its standard deviation over the $k=5$ folds. Every model with
$\bar{C} \ge \bar{C}^* - s^*/\sqrt{k}$ is statistically tied with the best, and among those the most explainable
model, the neuro-symbolic one, is chosen.

\subsection{Estimating the Month of Progression}
Most patients never progress, so the plain median of a Cox survival curve is usually undefined. The quantity
reported is the time to progression \emph{given} that progression occurs. Three families of estimators are
compared on held-out progressors:
\begin{itemize}
  \item the conditional median of the fused survival curve, the month at which $F(t)/F(\tau) = 1/2$ with
        $F = 1 - S$, for $\tau \in \{60, 120, 180\}$ months;
  \item log-normal and log-logistic AFT models \rc{aft} fitted on the training-fold progressors only, using the
        extended clinical design, the GNN risk, and the logic-engine risk;
  \item a constant baseline, the median progression month of the training-fold progressors.
\end{itemize}
The estimator with the lowest mean absolute error is used. Its interquartile range is reported with it, and the
estimate is mapped to one of four windows.

\subsection{Neuro-Symbolic Path Scoring and Top-K Selection}
A candidate explanation path is a walk over stored edges of the patient's fused graph:
patient $\rightarrow$ sample $\rightarrow$ altered gene $\rightarrow$ (protein partner) $\rightarrow$ pathway,
drug, or breast disease. Up to 60 candidates are enumerated per patient. Let $\bar{m}(p)$ be the mean
GNNExplainer importance of path $p$'s edges and $\operatorname{pct}(\cdot)$ its percentile among all candidate paths
in the cohort. Let $\mathcal{F}(p)$ be the logic-engine formulas, from the fold that held the patient out, that the
path \emph{grounds}: formulas that fire for the patient and have an atom that mentions the path's end node or
altered gene, or a relational PrimeKG atom $\mathit{Altered}(p,g) \wedge \mathit{InPathway}(g,w)$ whose gene $g$
lies on the path. With $S$ the 95th percentile of $|\sum_{f} w_f|$ over all rule-matched paths,
\begin{align}
  \text{neural}(p) &= \operatorname{pct}\big(\bar{m}(p)\big),\\
  \text{rule}(p) &= \min\!\Big(\big|\textstyle\sum_{f \in \mathcal{F}(p)} w_f\big| \,/\, S,\; 1\Big),\\
  \text{NS}(p) &= 0.6\,\text{neural}(p) + 0.4\,\text{rule}(p).
\end{align}
The sign of $\sum_f w_f$ gives the direction (faster or slower progression). The log-rank evidence of the end fact
(Section~\ref{sec:kgevidence}) is printed beside the path as cohort context, but no longer drives the score.
A path is \emph{prominent} when its NS score is in the top 20\% of all candidates in the cohort. A report shows
$K = \min(5, \text{number of prominent paths})$. At most two paths may pass through the same gene, and each gene
route and each end node appears once. When fewer than three paths are prominent, the best remaining paths are
added and marked as non-prominent.

\section{Report Groups}
Reports are grouped by the \emph{recorded} progression month, so a group never mixes true progressors with
patients who were not seen to progress (Table~\ref{tab:groups}). Only 12 of the 145 progressions happen within
six months, so the first group is ``within 12 months''.

\begin{table}[htbp]
  \centering
  \small\setlength{\tabcolsep}{4pt}
  \caption{Report groups by recorded progression-free survival.}
  \label{tab:groups}
  \begin{tabular}{llr}
    \toprule
    Group & Recorded progression & Patients \\
    \midrule
    A & $\le 12$ months & 32 \\
    B & 12--24 months & 30 \\
    C & 24--36 months & 27 \\
    D & $> 36$ months & 56 \\
    N & no progression recorded (censored) & 938 \\
    \bottomrule
  \end{tabular}
\end{table}

\section{Patient Report Design}
Each report is an A4 PDF of two to three pages, produced with ReportLab and Matplotlib. It contains:
\begin{enumerate}
  \item \textbf{Patient and tumour record}: demographics, PAM50 subtype, AJCC stage and TNM, histology, therapy
        history, aneuploidy, TMB, and MSI score.
  \item \textbf{Progression assessment}: the risk percentile and tier with a gauge, the probability of progression
        by 12, 24, 36, and 60 months, the estimated progression month and window, the recorded outcome with a
        tick or cross against the estimate, and the patient's survival curve against the cohort Kaplan--Meier curve.
  \item \textbf{Neuro-symbolic reasoning}: the top-K paths drawn as chains with their neural and symbolic score
        split, a table of the rule evidence behind each path, and the weighted logic formulas that fire for
        the patient, with their weight, direction, source (RL-proposed, mined, or single atom), support, and the
        Kaplan--Meier curve of the training patients for whom the strongest formula holds.
  \item \textbf{Molecular profile and knowledge-graph context}: altered genes with protein change, PrimeKG links,
        and GNN importance; PrimeKG drug links (research context only); and the five most similar patients by
        cosine similarity of their predicates, with their recorded outcomes.
\end{enumerate}
Each progression group also gets a summary page (group statistics, frequent path end-points, frequent rules,
common driver alterations, and a patient index) and a combined PDF. A cohort overview compares all models.

% ===========================================================================
\chapter{Experimental Setup and Results Analysis}\label{ch:results}
% ===========================================================================

\section{Experimental Setup}
\paragraph{Hardware and software.} All experiments ran on an Apple Silicon laptop (8 CPU cores, 8 GB RAM) without a
GPU. Software: Python 3.13, PyTorch 2.14, PyTorch Geometric 2.8, scikit-learn 1.9, lifelines 0.30, scikit-survival
0.28, RuleKit 2.1.24 on Java 22, and ReportLab 5.0. The full neuro-symbolic pipeline runs in about 50 minutes, and
the 1,083 reports build in about 5 minutes on six processes.

\paragraph{Protocol.} All evaluation is stratified five-fold cross-validation with seed 42. Within each training
fold, 15\% is held out for early stopping. Clinical standardisation is fitted on training folds only. Graph models
train for up to 30--40 epochs with patience 6--8, using Adam (learning rate $10^{-3}$, weight decay $10^{-4}$).

\paragraph{Metrics.} AUROC and AUPRC for the landmark label; Harrell's C (pooled, and the mean $\pm$ standard
deviation over folds); Uno's C truncated at 120 months; time-dependent AUC at 12, 36, and 60 months; and the
integrated Brier score up to 120 months. For the month of progression: mean and median absolute error, the share
of estimates within 12 months, window accuracy, and interquartile-range coverage.

\section{Milestone 1: Patient Graphs}
Graphs were built and validated for every patient. A raw-PFS binary GraphSAGE model on all 1,083 patients
reached only AUROC 0.542 and AUPRC 0.216, close to chance. Censoring makes ``no recorded progression'' an unreliable
negative label, which motivated the landmark endpoint.

\section{Milestone 2: PrimeKG Fusion and the 36-Month Benchmark}
Table~\ref{tab:36m} shows the first fused-graph benchmark on the 36-month landmark (498 patients, 89 events;
the prevalence, which is the AUPRC of a random ranking, is 0.179). The clinical logistic regression is the best
model. GCN is the best graph model and comes close without beating it. The gene multi-hot MLP is the weakest,
because the wide alteration vector adds noise. Figure~\ref{fig:roc36} shows the out-of-fold ROC curves of the
graph models.

\begin{table}[htbp]
  \centering
  \small\setlength{\tabcolsep}{4pt}
  \caption{36-month progression, five-fold mean $\pm$ standard deviation.}
  \label{tab:36m}
  \begin{tabular}{lcc}
    \toprule
    Model & AUROC & AUPRC \\
    \midrule
    Clinical logistic regression & $0.703 \pm 0.081$ & $0.463 \pm 0.133$ \\
    GCN & $0.680 \pm 0.076$ & $0.396 \pm 0.091$ \\
    GraphSAGE & $0.650 \pm 0.045$ & $0.379 \pm 0.091$ \\
    GATv2 & $0.631 \pm 0.046$ & $0.355 \pm 0.060$ \\
    MLP (clinical + gene multi-hot) & $0.562 \pm 0.098$ & $0.262 \pm 0.064$ \\
    \bottomrule
  \end{tabular}
\end{table}

\begin{figure}[htbp]
  \centering
  \includegraphics[width=0.62\textwidth]{fig_36m_roc.png}
  \caption{Out-of-fold ROC curves of the graph models on the 36-month endpoint.}
  \label{fig:roc36}
\end{figure}

\paragraph{Choice of horizon.} Figure~\ref{fig:horizon} sweeps the landmark from 12 to 60 months with the clinical
model. Short horizons have too few events (32 at 12 months). Long horizons drop too many censored patients (only
344 remain at 60 months). The 36-month horizon balances the two.

\begin{figure}[htbp]
  \centering
  \includegraphics[width=0.7\textwidth]{fig_horizon.png}
  \caption{Clinical AUROC as a function of the progression window.}
  \label{fig:horizon}
\end{figure}

\section{Concordance Training and the Architecture Zoo}
Table~\ref{tab:zoo} compares the seven architectures under binary cross-entropy and under the C-index objective
on the same folds. RGAT and GCN are the strongest graph models. GIN, which ignores relation types, is the weakest,
so typing the edges matters. The C-index objective does not improve discrimination over cross-entropy. No graph
model beats the clinical baseline (AUROC 0.703, C-index 0.682), and fold standard deviations of 0.05--0.15 exceed
most of the differences.

\begin{table}[htbp]
  \centering
  \small\setlength{\tabcolsep}{4pt}
  \caption{36-month endpoint: AUROC, AUPRC, and Harrell's C (five-fold means) under two objectives.}
  \label{tab:zoo}
  \resizebox{\textwidth}{!}{%
  \begin{tabular}{lccc|ccc}
    \toprule
    & \multicolumn{3}{c|}{Cross-entropy} & \multicolumn{3}{c}{C-index objective} \\
    Model & AUROC & AUPRC & C & AUROC & AUPRC & C \\
    \midrule
    RGAT & 0.694 & 0.412 & 0.672 & 0.678 & 0.389 & 0.658 \\
    GCN & 0.688 & 0.415 & 0.665 & 0.658 & 0.379 & 0.639 \\
    GraphSAGE & 0.679 & 0.406 & 0.658 & 0.660 & 0.393 & 0.643 \\
    GATv2 & 0.638 & 0.366 & 0.626 & 0.640 & 0.366 & 0.628 \\
    R-GCN & 0.627 & 0.352 & 0.616 & 0.602 & 0.343 & 0.597 \\
    Transformer & 0.613 & 0.354 & 0.607 & 0.627 & 0.376 & 0.617 \\
    GIN & 0.590 & 0.320 & 0.581 & 0.558 & 0.263 & 0.555 \\
    MLP & 0.570 & 0.274 & 0.562 & 0.566 & 0.263 & 0.560 \\
    \bottomrule
  \end{tabular}}
\end{table}

\section{Graph Structure Learning}
Figure~\ref{fig:gsl} compares the fixed two-hop graph with two GSL variants. The first variant used free edge
weights and let the scorer read the hop index. It closed the entire second hop (mean weight 0.24, and none of 5,745
hop-2 edges above 0.5), and GCN's AUPRC fell from 0.396 to 0.316. The hop-prior variant starts from the fixed graph
and gates only the third hop. Its run reported higher means (AUROC about 0.70 and AUPRC about 0.46 for all three
models), but within the fold standard deviations. That run's directory was not kept, so those bars are transcribed
from its report and cannot be recomputed.

\begin{figure}[htbp]
  \centering
  \includegraphics[width=\textwidth]{fig_gsl.png}
  \caption{Graph structure learning against the fixed two-hop graph on the 36-month endpoint.}
  \label{fig:gsl}
\end{figure}

\section{External Benchmark on GSE2034}
On the GSE2034 microarray cohort (286 patients, bone-relapse label, a 4,712-gene PrimeKG-aligned expression graph),
logistic regression reached AUROC $0.722 \pm 0.078$ and AUPRC 0.593. The random forest reached 0.682, GraphSAGE
0.582, GCN 0.568, and a majority baseline 0.500. As on TCGA, the expression graph did not lift the GNNs above a
linear baseline. A Cox model on principal components of the TCGA clinical block reached a test C-index of
$0.659 \pm 0.019$.

\section{Neuro-Symbolic Model}
\subsection{Cohort}
Figure~\ref{fig:pfshist} shows when the 145 recorded progressions occur. The median is 26 months, and progressions
continue past 10 years.

\begin{figure}[htbp]
  \centering
  \includegraphics[width=0.9\textwidth]{fig_pfs_hist.png}
  \caption{Recorded progression months, with the report group boundaries at 12, 24, and 36 months.}
  \label{fig:pfshist}
\end{figure}

\subsection{Discrimination and Model Selection}
Table~\ref{tab:ns} and Figure~\ref{fig:modelcmp} compare all models on the full follow-up, out-of-fold, on the same
five folds.

\begin{table}[htbp]
  \centering
  \small\setlength{\tabcolsep}{4pt}
  \caption{Full PFS follow-up (1,083 patients, 145 events). Harrell's C is the five-fold mean $\pm$ standard
  deviation. Uno's C is truncated at 120 months, and IBS is the integrated Brier score to 120 months (lower is better).}
  \label{tab:ns}
  \begin{tabular}{lccc}
    \toprule
    Model & Harrell's C & Uno's C & IBS \\
    \midrule
    Clinical Cox (19-field block) & $0.675 \pm 0.035$ & 0.616 & 0.134 \\
    Clinical Cox (extended) & $0.667 \pm 0.042$ & 0.657 & 0.135 \\
    \textbf{Neuro-symbolic Cox} (clin.\ + GNN + logic) & $\mathbf{0.664 \pm 0.021}$ & 0.616 & 0.134 \\
    Neuro-symbolic Cox (ext. clinical) & $0.662 \pm 0.042$ & 0.655 & 0.135 \\
    GNN + logic, no clinical & $0.642 \pm 0.056$ & 0.606 & 0.139 \\
    Survival MLN (logic alone) & $0.630 \pm 0.055$ & 0.617 & -- \\
    RL-structured MLN (logic engine alone) & $0.628 \pm 0.055$ & 0.615 & -- \\
    GraphSAGE (C-index objective) & $0.622 \pm 0.067$ & 0.573 & -- \\
    RuleKit survival rules & $0.619 \pm 0.048$ & 0.623 & -- \\
    GCN (Cox objective) & $0.612 \pm 0.044$ & 0.549 & -- \\
    GraphSAGE (Cox objective) & $0.603 \pm 0.062$ & 0.556 & -- \\
    R-GCN (Cox objective) & $0.586 \pm 0.041$ & 0.590 & -- \\
    Survival PSL & $0.604 \pm 0.059$ & 0.571 & -- \\
    RL rules & $0.563 \pm 0.046$ & 0.570 & -- \\
    Random survival forest & $0.554 \pm 0.023$ & 0.559 & 0.145 \\
    \bottomrule
  \end{tabular}
\end{table}

\begin{figure}[htbp]
  \centering
  \includegraphics[width=\textwidth]{fig_model_comparison.png}
  \caption{Fold-mean Harrell's C-index of every model. The report model is highlighted.}
  \label{fig:modelcmp}
\end{figure}

The best fold-mean C-index is 0.675 (clinical Cox), with one standard error of 0.016. Four models clear the bar of
0.659 and are statistically tied with it. Among them the neuro-symbolic Cox model is chosen because it is the one
whose prediction can be explained through paths and rules.

Two findings matter for interpretation:
\begin{enumerate}
  \item \textbf{Saturation.} GNNs (0.59--0.62), logic reasoners (0.56--0.63), forests (0.55), and Cox models
        (0.662--0.675) all stop near the same ceiling. Without the clinical block, the neural and logic risks reach
        0.642 together, better than either alone (0.622 and 0.628). Every Cox model that includes the clinical block
        reaches 0.662--0.675.
  \item \textbf{Explanation, not accuracy.} In every fold the fused model gives the GNN and logic terms hazard
        ratios of 0.94--1.09 per standard deviation with $p > 0.28$. Once the clinical block is present they add no
        independent prognostic information. The reports state this explicitly.
\end{enumerate}

\subsection{Symbolic Reasoners Compared}
Table~\ref{tab:logic} and Figure~\ref{fig:logic} compare the four logic reasoners with the RuleKit baseline on the
same folds.

\begin{table}[htbp]
  \centering
  \small\setlength{\tabcolsep}{4pt}
  \caption{Symbolic reasoners on the full follow-up: five-fold Harrell's C (mean $\pm$ sd), Uno's C, learned formulas
  per fold, and the share of formulas that reason over a PrimeKG relation.}
  \label{tab:logic}
  \resizebox{\textwidth}{!}{%
  \begin{tabular}{lcccc}
    \toprule
    Reasoner & Harrell's C & Uno's C & Formulas/fold & PrimeKG \\
    \midrule
    Survival MLN (Cox-MLN) & $0.630 \pm 0.055$ & 0.617 & 32.8 & 63\% \\
    \textbf{RL-structured survival MLN} (engine) & $\mathbf{0.628 \pm 0.055}$ & 0.615 & 35.2 & \textbf{69\%} \\
    RuleKit LR-Rules (baseline) & $0.619 \pm 0.048$ & 0.623 & 7.8 & 10\% \\
    Survival PSL (collective) & $0.604 \pm 0.059$ & 0.571 & 40.0 & 63\% \\
    RL rules (policy gradient) & $0.563 \pm 0.046$ & 0.570 & 25.6 & 91\% \\
    \bottomrule
  \end{tabular}}
\end{table}

\begin{figure}[htbp]
  \centering
  \includegraphics[width=0.9\textwidth]{fig_logic_comparison.png}
  \caption{Fold-mean Harrell's C of the symbolic reasoners. The engine is highlighted.}
  \label{fig:logic}
\end{figure}

\paragraph{What the comparison shows.}
\begin{enumerate}
  \item \textbf{Markov logic leads, but within noise.} Both MLN variants (0.630 and 0.628) are ahead of RuleKit
        (0.619), PSL (0.604), and RL rules alone (0.563). Their margin over RuleKit is smaller than one fold standard
        deviation, so the gain is not in discrimination. Under the one-standard-error rule (bar 0.605) MLN,
        RL-structured MLN, and RuleKit are tied, and the preference order selects the RL-structured MLN.
  \item \textbf{The gain is in the reasoning.} 69\% of the engine's formulas reason over PrimeKG relations, against
        10\% of RuleKit's rules. The formulas are weighted, combinable, and multi-hop. Across the folds the engine kept
        176 formulas: 51 single atoms, 81 mined clinical $\times$ PrimeKG pairs, and 44 multi-atom formulas proposed by
        reinforcement learning.
  \item \textbf{Collective inference did not help.} PSL learned a collective weight of about 0.01: in this cohort,
        similar patients carry no outcome signal beyond their own facts.
  \item \textbf{RL search overfits on its own.} The policy sampled 20,000--27,000 distinct rules per fold. Even with
        the replication filter, the rules it ranks highest do not transfer well (0.563). RL is useful as the structure
        learner of the MLN, where the Cox likelihood and its $\ell_1$ penalty decide which proposals carry weight.
\end{enumerate}

\paragraph{Learned formulas.} Clinical conditions still carry the largest single weights, for example
$\mathit{NStage}(p,\mathrm{N0}) \Rightarrow \mathit{Progress}(p)$ (weight $-0.32$ per SD) and
$\mathit{TStage}(p,\mathrm{T4})$ ($+0.13$). Typical PrimeKG formulas are a mined pair ($+0.15$, 43 patients),
\begin{multline*}
  \mathit{NStage}(p,\mathrm{N3}) \wedge \mathit{Altered}(p,g) \\
  \wedge \mathit{LinkedTo}(g,\text{breast ductal adenocarcinoma}) \Rightarrow \mathit{Progress}(p),
\end{multline*}
and an RL-proposed multi-hop formula ($-0.14$, 93 patients, a protective formula; Figure~\ref{fig:rulekm}),
\begin{multline*}
  \mathit{Altered}(p,g) \wedge \mathit{InPathway}(g,\text{VEGFA-VEGFR2}) \\
  \wedge \mathit{Altered}(p,g) \wedge \mathit{InPathway}(g,\text{Ion homeostasis}) \Rightarrow \mathit{Progress}(p).
\end{multline*}

\begin{figure}[htbp]
  \centering
  \includegraphics[width=0.6\textwidth]{fig_rule_km.png}
  \caption{Kaplan--Meier curve of the patients for whom the RL-proposed formula
  $\mathit{Altered}(p,g) \wedge \mathit{InPathway}(g,\text{VEGFA-VEGFR2}) \wedge \mathit{Altered}(p,g) \wedge
  \mathit{InPathway}(g,\text{Ion homeostasis})$ holds, against the cohort.}
  \label{fig:rulekm}
\end{figure}

\subsection{Estimating the Month of Progression}
Table~\ref{tab:timing} scores the estimators on the 145 held-out progressors. The conditional medians of the Cox
curve are poor, because proportional hazards fix the shape of the curve. The AFT models fitted on progressors are
best. The chosen log-logistic AFT model is within 12 months of the recorded month for 48\% of progressors (37\% for
the constant baseline) and in the recorded window for 33\% (19\% for the baseline, about 25\% by chance).
Figure~\ref{fig:timing} shows that estimates are compressed towards 15--45 months: early progressions are
overestimated and late ones underestimated. Agreement with the recorded window is 12\% in group A, 60\% in B, 37\%
in C, and 29\% in D. Several estimators were compared on the same progressors, so the chosen row is mildly
optimistic.

\begin{table}[htbp]
  \centering
  \small\setlength{\tabcolsep}{4pt}
  \caption{Estimating the month of progression on held-out progressors.}
  \label{tab:timing}
  \resizebox{\textwidth}{!}{%
  \begin{tabular}{lccccc}
    \toprule
    Estimator & MAE & Median AE & $\le 12$ mo & Window & IQR cov. \\
    \midrule
    Constant (training progressor median) & 22.9 & 16.0 & 37\% & 19\% & 49\% \\
    Conditional median, $\tau = 60$ & 23.3 & 18.2 & 33\% & 17\% & 37\% \\
    Conditional median, $\tau = 120$ & 28.0 & 24.6 & 23\% & 39\% & 41\% \\
    Conditional median, $\tau = 180$ & 59.4 & 63.1 & 7\% & 39\% & 32\% \\
    Log-normal AFT on progressors & 20.6 & 12.9 & 48\% & 34\% & 55\% \\
    \textbf{Log-logistic AFT on progressors} & \textbf{20.1} & \textbf{12.9} & \textbf{48\%} & \textbf{33\%} & \textbf{57\%} \\
    \bottomrule
  \end{tabular}}
\end{table}

\begin{figure}[htbp]
  \centering
  \includegraphics[width=0.6\textwidth]{fig_timing.png}
  \caption{Model-estimated against recorded progression month. Bars are the model's interquartile ranges, and the
  diagonal is perfect agreement.}
  \label{fig:timing}
\end{figure}

\subsection{Neuro-Symbolic Paths}
The explain stage enumerated 59,598 candidate paths for 1,033 patients in about 9 minutes. The other 50 patients
have no path, because none of their altered genes has a PrimeKG pathway, drug, or breast-disease link. The
prominence cutoff is an NS score of 0.523 (Figure~\ref{fig:nshist}). 27.5\% of all candidate paths, and 35\% of
the paths shown in reports, ground at least one learned formula. Most patients (801) receive $K=3$ paths, because
they have fewer than three prominent paths and are topped up; 62 patients have five prominent paths. A median of
three weighted formulas fires per patient, and 9.7\% of patients have none. Figure~\ref{fig:paths} shows the four
prominent paths of patient TCGA-BH-A1EV (group A, progression at 12.0 months). All four ground formulas that raise
the hazard. The strongest runs from the amplified gene NFASC to the pathway \emph{Neurofascin interactions}
($\sum w = +0.134$ over two formulas).

\begin{figure}[htbp]
  \centering
  \includegraphics[width=0.9\textwidth]{fig_ns_hist.png}
  \caption{Distribution of neuro-symbolic path scores, with the prominence cutoff at the cohort's 80th percentile.}
  \label{fig:nshist}
\end{figure}

\begin{figure}[htbp]
  \centering
  \includegraphics[width=\textwidth]{fig_paths_example.png}
  \caption{Top-K paths of patient TCGA-BH-A1EV. The bar on the right splits each NS score into its neural part
  (navy) and its rule part (red: the grounded formulas raise the hazard; blue: they lower it).}
  \label{fig:paths}
\end{figure}

\subsection{Patient Reports}
The system wrote 1,083 patient reports, five group summaries, four combined group PDFs, and a cohort overview.
Figures~\ref{fig:report1} to~\ref{fig:report3} show the full report of TCGA-BH-A1EV. Figure~\ref{fig:reportn}
shows the first page of a patient without recorded progression, whose outcome tile states the follow-up length and
that censoring is not proof of remission. Figure~\ref{fig:groupsummary} shows the summary page of group A.

\begin{figure}[htbp]
  \centering
  \includegraphics[width=0.8\textwidth]{fig_report_p1.png}
  \caption{Patient report, page 1: record, progression assessment, survival curve, and summary.}
  \label{fig:report1}
\end{figure}

\begin{figure}[htbp]
  \centering
  \includegraphics[width=0.8\textwidth]{fig_report_p2.png}
  \caption{Patient report, page 2: top-K knowledge-graph paths and the weighted logic formulas that fire.}
  \label{fig:report2}
\end{figure}

\begin{figure}[htbp]
  \centering
  \includegraphics[width=0.8\textwidth]{fig_report_p3.png}
  \caption{Patient report, page 3: molecular profile, PrimeKG drug links, similar patients, and reading notes.}
  \label{fig:report3}
\end{figure}

\begin{figure}[htbp]
  \centering
  \includegraphics[width=0.8\textwidth]{fig_report_n_p1.png}
  \caption{First page of a report in group N (no progression recorded).}
  \label{fig:reportn}
\end{figure}

\begin{figure}[htbp]
  \centering
  \includegraphics[width=0.8\textwidth]{fig_group_summary.png}
  \caption{Summary page of group A, progression within 12 months.}
  \label{fig:groupsummary}
\end{figure}

\section{Discussion}
Across five milestones, every attempt to raise discrimination with graph structure ended at the same ceiling:
relation-aware layers, concordance objectives, structure learning, and survival training on the full follow-up.
With 145 events, nodal status, stage, age, and subtype carry most of the prognostic signal, and somatic alterations
on a 500-gene panel add little beyond it. The PrimeKG facts are not uninformative. Several pathway facts separate
carriers with hazard ratios around 1.5 on the training folds, and the neural and logic risks together reach a
C-index of 0.642 without any clinical input. Their information, however, overlaps with the clinical variables. The
logic reasoners show the same pattern from the symbolic side. Moving from crisp IF--THEN rules to weighted,
first-order, knowledge-graph-grounded formulas changes \emph{what} the system can say about a patient much more
than \emph{how well} it ranks patients.

The neuro-symbolic system therefore earns its place through explanation. Each report states which PrimeKG facts the
graph model attended to, which weighted logic formulas ground those paths, and which formulas fire for the
patient. It
does so without claiming accuracy that the evaluation does not support.

% ===========================================================================
\chapter{Conclusion and Future Work}\label{ch:conclusion}
% ===========================================================================

\section{Conclusion}
We built a complete pipeline from raw TCGA tables to patient-specific knowledge graphs fused with PrimeKG,
benchmarked seven graph neural network architectures under three objectives, and turned the result into a
neuro-symbolic progression model with per-patient reports. Its symbolic engine replaces IF--THEN rule induction
with a survival Markov logic network whose first-order, knowledge-graph-grounded formulas are proposed by
reinforcement learning and weighted by the Cox partial likelihood. We compared it against probabilistic soft logic
with collective inference, RL rules alone, and RuleKit. The fused model reaches a concordance index of 0.664,
statistically tied with the best clinical model, and it estimates the month of progression to within a year for
about half of progressors. Every prediction is out-of-fold and leakage-checked.
Every one of the 1,083 reports explains its estimate with explicit knowledge-graph paths and weighted logic
formulas, and
places the patient in a progression-timeline group.

The central finding is that, on this cohort, knowledge-graph structure adds explanation rather than accuracy. We
report this result as it is, and the system is designed around it.

\section{Limitations}
\begin{itemize}
  \item Discrimination saturates at a C-index of about 0.67 with 145 events. The neuro-symbolic terms carry close to
        zero weight once clinical variables are present.
  \item The month-of-progression estimate is coarse. Late progressions (beyond five years) are hard to place.
  \item The logic engine adds reasoning, not accuracy: Markov logic matches RuleKit within noise, PSL's collective
        rule found no signal, and RL rules alone overfit their search.
  \item Paths explain the GNN, which is the weaker component. They are evidence, not causal claims.
  \item PrimeKG drug links include enzymes, carriers, and transporters, and are not treatment recommendations.
  \item One cohort, one seed, and no external validation of the reports. The fusion uses cross-validated stacking.
\end{itemize}

\section{Future Work}
\begin{itemize}
  \item \textbf{More modalities}: add expression, methylation, and protein data as node features, which the current
        graph excludes by design.
  \item \textbf{Pretrained knowledge-graph embeddings}: initialise PrimeKG node embeddings from TransE or a
        pretrained biomedical model instead of training them from scratch.
  \item \textbf{Richer RL reasoning}: let the policy walk multi-hop PrimeKG relation paths (as in MINERVA) rather than
        choose atoms, and train it jointly with the MLN weights (as in RNNLogic).
  \item \textbf{Joint MLN inference}: add formulas that link patients (shared pathway disruption) and run full MLN
        inference instead of the factorised per-patient score.
  \item \textbf{External validation}: apply the model and reports to METABRIC or GSE2990, whose relapse-free survival
        times are already extracted.
  \item \textbf{Clinician evaluation}: have oncologists rate the usefulness and correctness of the reports.
\end{itemize}

% ===========================================================================
\appendix
\chapter{Reproducibility}
% ===========================================================================

All code is in the project repository. From the repository root, the commands below build the Milestone~1 graphs, run the Milestone~2 benchmark, run the five neuro-symbolic stages (neural, symbolic, fusion, explain), write the reports and the notebook, draw the figures of this report, and run the 58 unit tests:

\begin{verbatim}
pip install -r requirements.txt
python scripts/build_patient_graphs.py
python scripts/train_and_evaluate_gnns.py \
    --endpoint pfs_36m --skip-sensitivity
python scripts/run_neurosymbolic_pipeline.py
python scripts/build_patient_reports.py --workers 4
python scripts/build_neurosymbolic_notebook.py
python report/make_figures.py
python -m unittest discover -s tests
\end{verbatim}

The neuro-symbolic pipeline caches each stage under \texttt{outputs/neurosymbolic/}. The reports are written to
\texttt{outputs/patient\_reports/}, in one folder per group. The method is documented in
\texttt{docs/neurosymbolic\_reports.md}, and the executed walkthrough is
\texttt{notebooks/neurosymbolic\_patient\_reports.ipynb}. RuleKit needs a Java runtime, version 11 or later.

\begin{table}[htbp]
  \centering
  \small\setlength{\tabcolsep}{4pt}
  \caption{Main source modules of the neuro-symbolic system.}
  \label{tab:modules}
  \begin{tabular}{>{\ttfamily}l p{7.4cm}}
    \toprule
    \normalfont Module & Responsibility \\
    \midrule
    src/dataset\_fusion.py & PrimeKG slice, neighbourhood expansion, fused PyG graphs \\
    src/gnn\_models.py & seven GNN architectures and the clinical residual \\
    src/survival.py & comparable pairs, C-index, Cox and C-index objectives \\
    src/neurosymbolic/features.py & clinical, genomic, and PrimeKG predicates \\
    src/neurosymbolic/rules.py & RuleKit baseline and readable rule text \\
    src/neurosymbolic/logic/atoms.py & first-order ground atoms, soft survival targets \\
    src/neurosymbolic/logic/mln.py & survival Markov logic network (Cox-MLN) \\
    src/neurosymbolic/logic/psl.py & survival probabilistic soft logic \\
    src/neurosymbolic/logic/rl.py & policy-gradient rule generator \\
    src/neurosymbolic/logic/runner.py & reasoner comparison and engine choice \\
    src/neurosymbolic/kg\_evidence.py & vectorised log-rank evidence for PrimeKG facts \\
    src/neurosymbolic/paths.py & path scoring and top-K selection \\
    src/neurosymbolic/timing.py & report groups and the progression-month estimators \\
    src/neurosymbolic/report\_pdf.py & patient reports, group summaries, cohort overview \\
    \bottomrule
  \end{tabular}
\end{table}

% ===========================================================================
% References: numbered list cited with \rc{key} (see the preamble).
% ===========================================================================
\renewcommand{\baselinestretch}{1.1}\normalsize
\chapter*{References}
\addcontentsline{toc}{chapter}{References}
\renewcommand{\labelenumi}{[\arabic{enumi}]}
\begin{enumerate}
\small
\item\label{ref:globocan} H.~Sung, J.~Ferlay, R.~L. Siegel, et al. Global cancer statistics 2020: GLOBOCAN estimates of incidence and mortality worldwide for 36 cancers in 185 countries. \emph{CA: A Cancer Journal for Clinicians}, 71(3):209--249, 2021.
\item\label{ref:pam50} J.~S. Parker, M.~Mullins, M.~C.~U. Cheang, et al. Supervised risk predictor of breast cancer based on intrinsic subtypes. \emph{Journal of Clinical Oncology}, 27(8):1160--1167, 2009.
\item\label{ref:primekg} P.~Chandak, K.~Huang, and M.~Zitnik. Building a knowledge graph to enable precision medicine. \emph{Scientific Data}, 10:67, 2023.
\item\label{ref:gcn} T.~N. Kipf and M.~Welling. Semi-supervised classification with graph convolutional networks. In \emph{International Conference on Learning Representations (ICLR)}, 2017.
\item\label{ref:graphsage} W.~L. Hamilton, R.~Ying, and J.~Leskovec. Inductive representation learning on large graphs. In \emph{Advances in Neural Information Processing Systems (NeurIPS)}, 2017.
\item\label{ref:gnnexplainer} R.~Ying, D.~Bourgeois, J.~You, M.~Zitnik, and J.~Leskovec. GNNExplainer: Generating explanations for graph neural networks. In \emph{Advances in Neural Information Processing Systems (NeurIPS)}, 2019.
\item\label{ref:neurosymbolic} A.~d'Avila Garcez and L.~C. Lamb. Neurosymbolic AI: the 3rd wave. \emph{Artificial Intelligence Review}, 56:12387--12406, 2023.
\item\label{ref:cbioportal} E.~Cerami, J.~Gao, U.~Dogrusoz, et al. The cBio cancer genomics portal: an open platform for exploring multidimensional cancer genomics data. \emph{Cancer Discovery}, 2(5):401--404, 2012.
\item\label{ref:cbioportal2} J.~Gao, B.~A. Aksoy, U.~Dogrusoz, et al. Integrative analysis of complex cancer genomics and clinical profiles using the cBioPortal. \emph{Science Signaling}, 6(269):pl1, 2013.
\item\label{ref:rulekit} A.~Gudy\'s, M.~Sikora, and \L.~Wr\'obel. RuleKit: A comprehensive suite for rule-based learning. \emph{Knowledge-Based Systems}, 194:105480, 2020.
\item\label{ref:hoadley} K.~A. Hoadley, C.~Yau, T.~Hinoue, et al. Cell-of-origin patterns dominate the molecular classification of 10,000 tumors from 33 types of cancer. \emph{Cell}, 173(2):291--304, 2018.
\item\label{ref:tcgacdr} J.~Liu, T.~Lichtenberg, K.~A. Hoadley, et al. An integrated TCGA pan-cancer clinical data resource to drive high-quality survival outcome analytics. \emph{Cell}, 173(2):400--416, 2018.
\item\label{ref:gistic} C.~H. Mermel, S.~E. Schumacher, B.~Hill, et al. GISTIC2.0 facilitates sensitive and confident localization of the targets of focal somatic copy-number alteration in human cancers. \emph{Genome Biology}, 12:R41, 2011.
\item\label{ref:reactome} B.~Jassal, L.~Matthews, G.~Viteri, et al. The Reactome pathway knowledgebase. \emph{Nucleic Acids Research}, 48(D1):D498--D503, 2020.
\item\label{ref:drugbank} D.~S. Wishart, Y.~D. Feunang, A.~C. Guo, et al. DrugBank 5.0: a major update to the DrugBank database for 2018. \emph{Nucleic Acids Research}, 46(D1):D1074--D1082, 2018.
\item\label{ref:gatv2} S.~Brody, U.~Alon, and E.~Yahav. How attentive are graph attention networks? In \emph{International Conference on Learning Representations (ICLR)}, 2022.
\item\label{ref:rgcn} M.~Schlichtkrull, T.~N. Kipf, P.~Bloem, R.~van den Berg, I.~Titov, and M.~Welling. Modeling relational data with graph convolutional networks. In \emph{The Semantic Web (ESWC)}, pages 593--607, 2018.
\item\label{ref:rgat} D.~Busbridge, D.~Sherburn, P.~Cavallo, and N.~Y. Hammerla. Relational graph attention networks. \emph{arXiv preprint arXiv:1904.05811}, 2019.
\item\label{ref:gin} K.~Xu, W.~Hu, J.~Leskovec, and S.~Jegelka. How powerful are graph neural networks? In \emph{International Conference on Learning Representations (ICLR)}, 2019.
\item\label{ref:transformerconv} Y.~Shi, Z.~Huang, S.~Feng, H.~Zhong, W.~Wang, and Y.~Sun. Masked label prediction: Unified message passing model for semi-supervised classification. In \emph{Proceedings of IJCAI}, pages 1548--1554, 2021.
\item\label{ref:pyg} M.~Fey and J.~E. Lenssen. Fast graph representation learning with PyTorch Geometric. In \emph{ICLR Workshop on Representation Learning on Graphs and Manifolds}, 2019.
\item\label{ref:gslsurvey} Y.~Zhu, W.~Xu, J.~Zhang, et al. A survey on graph structure learning: Progress and opportunities. \emph{arXiv preprint arXiv:2103.03036}, 2021.
\item\label{ref:lds} L.~Franceschi, M.~Niepert, M.~Pontil, and X.~He. Learning discrete structures for graph neural networks. In \emph{International Conference on Machine Learning (ICML)}, pages 1972--1982, 2019.
\item\label{ref:km} E.~L. Kaplan and P.~Meier. Nonparametric estimation from incomplete observations. \emph{Journal of the American Statistical Association}, 53(282):457--481, 1958.
\item\label{ref:logrank} N.~Mantel. Evaluation of survival data and two new rank order statistics arising in its consideration. \emph{Cancer Chemotherapy Reports}, 50(3):163--170, 1966.
\item\label{ref:cox} D.~R. Cox. Regression models and life-tables. \emph{Journal of the Royal Statistical Society: Series B}, 34(2):187--220, 1972.
\item\label{ref:aft} L.~J. Wei. The accelerated failure time model: a useful alternative to the Cox regression model in survival analysis. \emph{Statistics in Medicine}, 11(14--15):1871--1879, 1992.
\item\label{ref:rsf} H.~Ishwaran, U.~B. Kogalur, E.~H. Blackstone, and M.~S. Lauer. Random survival forests. \emph{The Annals of Applied Statistics}, 2(3):841--860, 2008.
\item\label{ref:deepsurv} J.~L. Katzman, U.~Shaham, A.~Cloninger, J.~Bates, T.~Jiang, and Y.~Kluger. DeepSurv: personalized treatment recommender system using a Cox proportional hazards deep neural network. \emph{BMC Medical Research Methodology}, 18:24, 2018.
\item\label{ref:harrell} F.~E. Harrell, R.~M. Califf, D.~B. Pryor, K.~L. Lee, and R.~A. Rosati. Evaluating the yield of medical tests. \emph{JAMA}, 247(18):2543--2546, 1982.
\item\label{ref:uno} H.~Uno, T.~Cai, M.~J. Pencina, R.~B. D'Agostino, and L.~J. Wei. On the C-statistics for evaluating overall adequacy of risk prediction procedures with censored survival data. \emph{Statistics in Medicine}, 30(10):1105--1117, 2011.
\item\label{ref:lifelines} C.~Davidson-Pilon. lifelines: survival analysis in Python. \emph{Journal of Open Source Software}, 4(40):1317, 2019.
\item\label{ref:sksurv} S.~P\"olsterl. scikit-survival: A library for time-to-event analysis built on top of scikit-learn. \emph{Journal of Machine Learning Research}, 21(212):1--6, 2020.
\item\label{ref:lrrules} \L.~Wr\'obel, A.~Gudy\'s, and M.~Sikora. Learning rule sets from survival data. \emph{BMC Bioinformatics}, 18:285, 2017.
\item\label{ref:gse2034} Y.~Wang, J.~G.~M. Klijn, Y.~Zhang, et al. Gene-expression profiles to predict distant metastasis of lymph-node-negative primary breast cancer. \emph{The Lancet}, 365(9460):671--679, 2005.
\item\label{ref:gse2990} C.~Sotiriou, P.~Wirapati, S.~Loi, et al. Gene expression profiling in breast cancer: understanding the molecular basis of histologic grade to improve prognosis. \emph{Journal of the National Cancer Institute}, 98(4):262--272, 2006.
\item\label{ref:stacking} D.~H. Wolpert. Stacked generalization. \emph{Neural Networks}, 5(2):241--259, 1992.
\item\label{ref:mln} M.~Richardson and P.~Domingos. Markov logic networks. \emph{Machine Learning}, 62(1--2):107--136, 2006.
\item\label{ref:mlnstructure} S.~Kok and P.~Domingos. Learning the structure of Markov logic networks. In \emph{International Conference on Machine Learning (ICML)}, pages 441--448, 2005.
\item\label{ref:psl} S.~H. Bach, M.~Broecheler, B.~Huang, and L.~Getoor. Hinge-loss Markov random fields and probabilistic soft logic. \emph{Journal of Machine Learning Research}, 18(109):1--67, 2017.
\item\label{ref:reinforce} R.~J. Williams. Simple statistical gradient-following algorithms for connectionist reinforcement learning. \emph{Machine Learning}, 8(3--4):229--256, 1992.
\item\label{ref:deeppath} W.~Xiong, T.~Hoang, and W.~Y. Wang. DeepPath: A reinforcement learning method for knowledge graph reasoning. In \emph{Proceedings of EMNLP}, pages 564--573, 2017.
\item\label{ref:minerva} R.~Das, S.~Dhuliawala, M.~Zaheer, et al. Go for a walk and arrive at the answer: Reasoning over paths in knowledge bases using reinforcement learning. In \emph{International Conference on Learning Representations (ICLR)}, 2018.
\item\label{ref:rnnlogic} M.~Qu, J.~Chen, L.-P. Xhonneux, Y.~Bengio, and J.~Tang. RNNLogic: Learning logic rules for reasoning on knowledge graphs. In \emph{International Conference on Learning Representations (ICLR)}, 2021.
\item\label{ref:gru} K.~Cho, B.~van Merri\"enboer, C.~Gulcehre, et al. Learning phrase representations using RNN encoder--decoder for statistical machine translation. In \emph{Proceedings of EMNLP}, pages 1724--1734, 2014.
\item\label{ref:elasticnet} H.~Zou and T.~Hastie. Regularization and variable selection via the elastic net. \emph{Journal of the Royal Statistical Society: Series B}, 67(2):301--320, 2005.
\item\label{ref:coxnet} N.~Simon, J.~Friedman, T.~Hastie, and R.~Tibshirani. Regularization paths for Cox's proportional hazards model via coordinate descent. \emph{Journal of Statistical Software}, 39(5):1--13, 2011.
\item\label{ref:cart} L.~Breiman, J.~H. Friedman, R.~A. Olshen, and C.~J. Stone. \emph{Classification and Regression Trees}. Wadsworth, Belmont, CA, 1984.
\end{enumerate}

\end{document}
